\documentclass[aspectratio=169,10pt]{beamer}

% Shared Beamer setup for Fall 2026 course slides.
% Clean Cal Poly-inspired theme: no custom class files or uncommon packages.
\mode<presentation>{
  \usetheme{default}
  \usefonttheme{professionalfonts}
}

\definecolor{CalPolyGreen}{HTML}{154734}
\definecolor{CalPolyGold}{HTML}{BD8B13}
\definecolor{CalPolySage}{HTML}{7FA28F}
\definecolor{CalPolyMint}{HTML}{DDF1E7}
\definecolor{CalPolySand}{HTML}{A29F86}
\definecolor{CalPolyBeige}{HTML}{EFEEDE}
\definecolor{DeckInk}{HTML}{1F2933}
\definecolor{DeckMuted}{HTML}{5F6B7A}
\definecolor{DeckSoft}{HTML}{F7F7F5}

\setbeamersize{text margin left=0.55in,text margin right=0.55in}
\setbeamercolor{normal text}{fg=DeckInk,bg=white}
\setbeamercolor{structure}{fg=CalPolyGreen}
\setbeamercolor{title}{fg=white}
\setbeamercolor{frametitle}{fg=CalPolyGreen,bg=white}
\setbeamercolor{block title}{bg=CalPolySage,fg=white}
\setbeamercolor{block body}{bg=CalPolyMint,fg=black}
\setbeamercolor{alerted text}{fg=CalPolyGold}

\setbeamerfont{title}{size=\Huge,series=\bfseries}
\setbeamerfont{frametitle}{size=\Large,series=\bfseries}
\setbeamerfont{block title}{size=\normalsize,series=\bfseries}

\setbeamertemplate{navigation symbols}{}
\setbeamertemplate{itemize item}{\color{CalPolyGold}$\blacktriangleright$}
\setbeamertemplate{itemize subitem}{\color{CalPolyGreen}$\bullet$}
\setbeamertemplate{footline}{
  \hfill{\color{DeckMuted}\scriptsize\insertframenumber/\inserttotalframenumber}
  \hspace{0.18in}\vspace{0.12in}
}
\setbeamertemplate{frametitle}{
  \vspace{0.18in}
  \hspace{0.02in}\insertframetitle\par
  \vspace{0.08in}
  \noindent\makebox[\linewidth][l]{%
    {\color{CalPolyGold}\rule{0.55in}{2pt}}%
    {\color{CalPolyGreen}\rule{\dimexpr\linewidth-0.55in\relax}{2pt}}%
  }
}

\newcommand{\CourseNumber}{Course}
\newcommand{\CourseTitle}{Course Title}
\newcommand{\LectureNumber}{0}
\newcommand{\LectureTitle}{Lecture Title}
\newcommand{\LectureDate}{Date}
\newcommand{\CourseUnit}{Unit}

\newcommand{\coursenumber}[1]{\renewcommand{\CourseNumber}{#1}}
\newcommand{\coursetitle}[1]{\renewcommand{\CourseTitle}{#1}}
\newcommand{\lecturenumber}[1]{\renewcommand{\LectureNumber}{#1}}
\newcommand{\lecturetitle}[1]{\renewcommand{\LectureTitle}{#1}}
\newcommand{\lecturedate}[1]{\renewcommand{\LectureDate}{#1}}
\newcommand{\courseunit}[1]{\renewcommand{\CourseUnit}{#1}}

\author{Kirk Duran}
\institute{California Polytechnic State University, San Luis Obispo}
\date{\LectureDate}

\newcommand{\makedecktitle}{
  \begingroup
  \setbeamercolor{background canvas}{bg=CalPolyGreen}
  \begin{frame}[plain]
    \vfill
    \begin{center}
      {\usebeamerfont{title}\color{white}\LectureTitle\par}
      \vspace{0.18in}
      {\Large\color{white}Lecture \LectureNumber\par}
    \end{center}
    \vfill
    \noindent{\color{white}\rule{\linewidth}{1.2pt}}\par
    \vspace{0.08in}
    \noindent\colorbox{CalPolyGold}{\makebox[0.24\linewidth][c]{\color{white}\Large\CourseNumber}}
    \hspace{0.12in}{\color{white}\Large Kirk Duran}
  \end{frame}
  \endgroup
}

\newenvironment{keyidea}{\begin{block}{Key idea}}{\end{block}}
\newenvironment{activity}{\begin{block}{In class}}{\end{block}}
\newenvironment{cpdefinition}{\begin{block}{Definition}}{\end{block}}
\newenvironment{exampleblockcp}{
  \begingroup
  \setbeamercolor{block title}{bg=CalPolySand,fg=white}
  \setbeamercolor{block body}{bg=CalPolyBeige,fg=black}
  \begin{block}{Example}
}{
  \end{block}
  \endgroup
}

\newcommand{\imageplaceholder}[2][0.3in]{%
  \noindent\makebox[\linewidth][c]{%
    \fcolorbox{CalPolySage}{CalPolyBeige}{%
      \parbox[c][#1][c]{0.86\linewidth}{%
        \centering
        {\color{DeckMuted}\scriptsize Image placeholder: }%
        {\color{DeckInk}\scriptsize\ttfamily\detokenize{#2}}%
      }%
    }%
  }%
}

\newcommand{\sectionframe}[1]{
  \begingroup
  \setbeamercolor{background canvas}{bg=CalPolyGreen}
  \begin{frame}[plain]
    \vfill
    {\color{white}\LARGE\bfseries #1\par}
    \vspace{0.18in}
    {\color{CalPolyGold}\rule{0.22\paperwidth}{3pt}}
    {\color{white}\rule{0.62\paperwidth}{3pt}}
    \vfill
  \end{frame}
  \endgroup
}

\newcommand{\placeholdernote}{
  \begin{block}{Placeholder}
    Replace these bullets with the final lecture material, examples, and in-class activities.
  \end{block}
}


\usepackage{tikz}

\graphicspath{{../assets/lecture-16/}}

% If an image file is not yet available, skip the visual slot.
\newcommand{\lectureimage}[2][1.75in]{%
  \IfFileExists{../assets/lecture-16/#2}{%
    \includegraphics[width=\linewidth,height=#1,keepaspectratio]{#2}%
  }{}%
}

% Explicit placeholder for visuals/examples that will be added later.
\newcommand{\lectureplaceholder}[2][2.35in]{%
  \begin{center}
    \fbox{%
      \begin{minipage}[c][#1][c]{0.90\linewidth}
        \centering
        \small\textit{#2}
      \end{minipage}%
    }
  \end{center}%
}

\setbeamersize{text margin left=0.42in,text margin right=0.42in}

\coursenumber{CS 4880}
\coursetitle{Artificial Intelligence}
\lecturenumber{15}
\lecturetitle{First-Order Logic and Knowledge Representation}
\lecturedate{October 2, 2026}
\courseunit{Logic and Knowledge Representation}

\begin{document}

% Estimated pacing: 50 minutes
% Why propositional logic is not enough: 7 min
% FOL vocabulary and sentence structure: 14 min
% Quantifiers and translation: 14 min
% Semantics and knowledge representation: 11 min
% Bridge to first-order inference: 4 min

% -----------------------------------------------------------------------------
% Slide 1
% -----------------------------------------------------------------------------
\makedecktitle

\section{Why First-Order Logic?}

% -----------------------------------------------------------------------------
% Slide 2
% -----------------------------------------------------------------------------
\begin{frame}[shrink=10]{Propositional Logic Has a Scaling Problem of Meaning}
  \begin{columns}[T,onlytextwidth]
    \begin{column}{0.54\textwidth}
      Propositional logic can encode individual facts and rules very cleanly.

      \begin{itemize}
        \item But each distinct object usually needs its own proposition.
        \item The same rule may have to be repeated for Alice, Bob, Carol, and every future object.
        \item The representation becomes large even when the underlying idea is simple.
      \end{itemize}
    \end{column}
    \begin{column}{0.40\textwidth}
      \lectureimage[2.65in]{slide2.png}
    \end{column}
  \end{columns}
  \begin{keyidea}
    First-order logic lets one rule describe an entire class of objects.
  \end{keyidea}
\end{frame}

% -----------------------------------------------------------------------------
% Slide 3
% -----------------------------------------------------------------------------
\begin{frame}[shrink=12]{The Rule We Actually Want to Say}
  Instead of writing one rule for every person:

  \[
    StudentAlice \rightarrow HasIDAlice
  \]
  \[
    StudentBob \rightarrow HasIDBob
  \]
  \[
    StudentCarol \rightarrow HasIDCarol
  \]

  we want one general rule:
  \[
    \forall x\; Student(x) \rightarrow HasID(x).
  \]

  \begin{keyidea}
    The variable $x$ lets one sentence stand for arbitrarily many individuals.
  \end{keyidea}
\end{frame}

% -----------------------------------------------------------------------------
% Slide 4
% -----------------------------------------------------------------------------
\begin{frame}[shrink=10]{Logical Agents Need More Than True/False Labels}
  \begin{columns}[T,onlytextwidth]
    \begin{column}{0.54\textwidth}
      An intelligent agent reasons about \emph{objects} in a world.

      \begin{itemize}
        \item Objects have properties.
        \item Objects participate in relationships.
        \item Objects belong to categories.
        \item The same relationships recur across many objects.
      \end{itemize}

      The knowledge base should preserve this structure rather than flattening everything into unrelated Boolean symbols.
    \end{column}
    \begin{column}{0.40\textwidth}
      \lectureimage[2.70in]{slide4.png}
    \end{column}
  \end{columns}
  \begin{keyidea}
    Knowledge representation is the choice of structure used to encode what the agent knows.
  \end{keyidea}
\end{frame}

% -----------------------------------------------------------------------------
% Slide 5
% -----------------------------------------------------------------------------
\begin{frame}[shrink=12]{From GOFAI to Knowledge Representation}
  \begin{columns}[T,onlytextwidth]
    \begin{column}{0.47\textwidth}
      \begin{block}{Symbolic AI}
        Classical AI treated intelligent behavior as manipulation of explicit symbols and structured knowledge.
      \end{block}
    \end{column}
    \begin{column}{0.47\textwidth}
      \begin{block}{Knowledge representation}
        Researchers developed logical languages, semantic networks, frames, rules, and later ontologies to organize facts about domains.
      \end{block}
    \end{column}
  \end{columns}

  \lectureimage[1.35in]{slide5.png}

  \begin{keyidea}
    The representation is not decoration: it determines what the system can express and what it can infer.
  \end{keyidea}
\end{frame}

% -----------------------------------------------------------------------------
% Slide 6
% -----------------------------------------------------------------------------
\begin{frame}[shrink=10]{A Running Example: A Campus Delivery Robot}
  \begin{columns}[T,onlytextwidth]
    \begin{column}{0.54\textwidth}
      Our agent operates in a campus world containing
      \begin{itemize}
        \item people;
        \item packages;
        \item rooms and buildings;
        \item doors and hallways;
        \item charging stations;
        \item other robots.
      \end{itemize}

      The same structures recur constantly: \emph{located in}, \emph{owns}, \emph{carries}, \emph{connected to}, and \emph{accessible to}.
    \end{column}
    \begin{column}{0.40\textwidth}
      \lectureimage[2.75in]{slide6.png}
    \end{column}
  \end{columns}
  \begin{keyidea}
    A good representation captures repeated structure instead of inventing a new proposition for every situation.
  \end{keyidea}
\end{frame}

\section{The Language of First-Order Logic}

% -----------------------------------------------------------------------------
% Slide 7
% -----------------------------------------------------------------------------
\begin{frame}[shrink=10]{First-Order Logic: The Big Upgrade}
  \begin{columns}[T,onlytextwidth]
    \begin{column}{0.31\textwidth}
      \begin{block}{Objects}
        Individual things in the domain.

        Alice, Robot7, Room214, PackageA
      \end{block}
    \end{column}
    \begin{column}{0.31\textwidth}
      \begin{block}{Properties / relations}
        Predicates describe objects and how they are connected.

        $Student(Alice)$

        $In(Robot7,Room214)$
      \end{block}
    \end{column}
    \begin{column}{0.31\textwidth}
      \begin{block}{Quantifiers}
        Variables let a sentence range over many objects.

        $\forall x$ and $\exists x$
      \end{block}
    \end{column}
  \end{columns}

  \lectureimage[1.20in]{slide7.png}

  \begin{keyidea}
    FOL represents a world in terms of entities and structured statements about those entities.
  \end{keyidea}
\end{frame}

% -----------------------------------------------------------------------------
% Slide 8
% -----------------------------------------------------------------------------
\begin{frame}[shrink=10]{Vocabulary: Constants and Variables}
  \begin{columns}[T,onlytextwidth]
    \begin{column}{0.54\textwidth}
      \begin{itemize}
        \item A \textbf{constant} names a particular object in the domain.
        \item A \textbf{variable} stands for an unspecified object.
        \item Variables become powerful when combined with quantifiers.
      \end{itemize}

      Examples:
      \[
        Alice,\; Robot7,\; Room214
      \]
      \[
        x,\; y,\; z
      \]
    \end{column}
    \begin{column}{0.40\textwidth}
      \lectureimage[2.50in]{slide8.png}
    \end{column}
  \end{columns}
  \begin{keyidea}
    Constants point to specific objects; variables let rules talk about objects generically.
  \end{keyidea}
\end{frame}

% -----------------------------------------------------------------------------
% Slide 9
% -----------------------------------------------------------------------------
\begin{frame}[shrink=12]{Constants Name Particular Objects}
  Examples from our campus domain:
  \[
    Alice
  \]
  \[
    Robot7
  \]
  \[
    Room214
  \]
  \[
    PackageA
  \]

  \begin{keyidea}
    A constant is a symbol interpreted as one particular object in the domain.
  \end{keyidea}
\end{frame}

% -----------------------------------------------------------------------------
% Slide 10
% -----------------------------------------------------------------------------
\begin{frame}[shrink=10]{Predicates Describe Properties}
  \begin{columns}[T,onlytextwidth]
    \begin{column}{0.54\textwidth}
      A one-place predicate behaves like a property test:
      \[
        Student(Alice)
      \]
      \[
        Robot(Robot7)
      \]
      \[
        Locked(Door3)
      \]
      \[
        Fragile(PackageA)
      \]
    \end{column}
    \begin{column}{0.40\textwidth}
      \lectureimage[2.55in]{slide10.png}
    \end{column}
  \end{columns}
  \begin{keyidea}
    An atomic sentence is true or false depending on the object and the interpretation.
  \end{keyidea}
\end{frame}

% -----------------------------------------------------------------------------
% Slide 11
% -----------------------------------------------------------------------------
\begin{frame}[shrink=10]{Relations Connect Objects}
  \begin{columns}[T,onlytextwidth]
    \begin{column}{0.54\textwidth}
      Predicates can take multiple arguments:
      \[
        In(Robot7, Room214)
      \]
      \[
        Carries(Robot7, PackageA)
      \]
      \[
        Teaches(ProfLee, AI)
      \]
      \[
        Connected(Room214, Hallway2)
      \]
    \end{column}
    \begin{column}{0.40\textwidth}
      \lectureimage[2.60in]{slide11.png}
    \end{column}
  \end{columns}
  \begin{keyidea}
    Relations preserve structure that propositional symbols hide.
  \end{keyidea}
\end{frame}

% -----------------------------------------------------------------------------
% Slide 12
% -----------------------------------------------------------------------------
\begin{frame}[shrink=10]{Argument Order Matters}
  \begin{columns}[T,onlytextwidth]
    \begin{column}{0.54\textwidth}
      \begin{itemize}
        \item $Parent(Alice,Bob)$ is generally not the same as $Parent(Bob,Alice)$.
        \item $In(Robot7,Room214)$ differs from $In(Room214,Robot7)$.
        \item A relation has meaning tied to both its name and the positions of its arguments.
      \end{itemize}
    \end{column}
    \begin{column}{0.40\textwidth}
      \lectureimage[2.50in]{slide12.png}
    \end{column}
  \end{columns}
  \begin{keyidea}
    Predicate arity and argument order are part of the representation.
  \end{keyidea}
\end{frame}

% -----------------------------------------------------------------------------
% Slide 13
% -----------------------------------------------------------------------------
\begin{frame}[shrink=10]{Functions Denote Objects}
  A function maps one or more objects to another object:
  \[
    OwnerOf(PackageA)
  \]
  \[
    SupervisorOf(Alice)
  \]
  \[
    NextRoom(Room214)
  \]
  \[
    MotherOf(x)
  \]

  \begin{keyidea}
    Predicates evaluate to truth values; functions denote objects.
  \end{keyidea}
\end{frame}

% -----------------------------------------------------------------------------
% Slide 14
% -----------------------------------------------------------------------------
\begin{frame}[shrink=10]{Predicates vs. Functions}
  \begin{columns}[T,onlytextwidth]
    \begin{column}{0.47\textwidth}
      \begin{block}{Predicate}
        Returns a truth value.

        \[
          In(Robot7, Room214)
        \]

        true or false
      \end{block}
    \end{column}
    \begin{column}{0.47\textwidth}
      \begin{block}{Function}
        Denotes an object.

        \[
          OwnerOf(PackageA)
        \]

        some person
      \end{block}
    \end{column}
  \end{columns}

  \begin{keyidea}
    Ask ``is this statement true?'' $\rightarrow$ predicate. Ask ``which object does this term refer to?'' $\rightarrow$ function.
  \end{keyidea}
\end{frame}

% -----------------------------------------------------------------------------
% Slide 15
% -----------------------------------------------------------------------------
\begin{frame}[shrink=10]{Terms Build References to Objects}
  \begin{columns}[T,onlytextwidth]
    \begin{column}{0.54\textwidth}
      \begin{itemize}
        \item A constant is a term.
        \item A variable is a term.
        \item A function applied to terms produces another term.
        \item Terms can therefore be nested.
      \end{itemize}

      Example:
      \[
        SupervisorOf(OwnerOf(PackageA))
      \]
    \end{column}
    \begin{column}{0.40\textwidth}
      \lectureimage[2.50in]{slide15.png}
    \end{column}
  \end{columns}
  \begin{keyidea}
    Terms refer to objects; predicates applied to terms produce sentences.
  \end{keyidea}
\end{frame}

% -----------------------------------------------------------------------------
% Slide 16
% -----------------------------------------------------------------------------
\begin{frame}[shrink=12]{Atomic Sentences Are Predicates Applied to Terms}
  Examples:
  \[
    Student(Alice)
  \]
  \[
    In(Robot7, NextRoom(Room214))
  \]
  \[
    Knows(Alice, SupervisorOf(Bob))
  \]

  \begin{keyidea}
    FOL combines a small vocabulary into richly structured statements.
  \end{keyidea}
\end{frame}

% -----------------------------------------------------------------------------
% Slide 17
% -----------------------------------------------------------------------------
\begin{frame}[shrink=12]{Compound FOL Sentences Reuse Familiar Connectives}
  The propositional connectives still work:
  \[
    Student(Alice) \land Enrolled(Alice,AI)
  \]
  \[
    Locked(Door3) \rightarrow \neg Accessible(Room214)
  \]
  \[
    Robot(Robot7) \land Carries(Robot7,PackageA)
  \]

  \begin{keyidea}
    First-order logic extends propositional logic; it does not replace AND, OR, NOT, and implication.
  \end{keyidea}
\end{frame}

\section{Quantifiers and Scope}

% -----------------------------------------------------------------------------
% Slide 18
% -----------------------------------------------------------------------------
\begin{frame}[shrink=10]{The Universal Quantifier: $\forall$}
  \begin{columns}[T,onlytextwidth]
    \begin{column}{0.54\textwidth}
      \begin{itemize}
        \item $\forall x$ means ``for every object $x$ in the domain.''
        \item It turns a pattern involving $x$ into a claim about all objects.
        \item Universal rules are one of the main reasons FOL is compact.
      \end{itemize}
    \end{column}
    \begin{column}{0.40\textwidth}
      \lectureimage[2.55in]{slide18.png}
    \end{column}
  \end{columns}
  \begin{keyidea}
    Use $\forall$ when the statement is intended to apply to every relevant object.
  \end{keyidea}
\end{frame}

% -----------------------------------------------------------------------------
% Slide 19
% -----------------------------------------------------------------------------
\begin{frame}[shrink=12]{Universal Rule Example}
  ``Every student has an ID.''
  \[
    \forall x\; Student(x) \rightarrow HasID(x)
  \]

  Suppose we also know
  \[
    Student(Alice)
  \]
  and
  \[
    Student(Bob).
  \]

  \begin{keyidea}
    One rule represents the entire category instead of repeating the rule for each person.
  \end{keyidea}
\end{frame}

% -----------------------------------------------------------------------------
% Slide 20
% -----------------------------------------------------------------------------
\begin{frame}[shrink=10]{The Existential Quantifier: $\exists$}
  \begin{columns}[T,onlytextwidth]
    \begin{column}{0.54\textwidth}
      \begin{itemize}
        \item $\exists x$ means ``there exists at least one object $x$.''
        \item The sentence does not need to identify which object satisfies the claim.
        \item Existential statements assert the existence of a witness.
      \end{itemize}
    \end{column}
    \begin{column}{0.40\textwidth}
      \lectureimage[2.50in]{slide20.png}
    \end{column}
  \end{columns}
  \begin{keyidea}
    Use $\exists$ when at least one object must satisfy the condition.
  \end{keyidea}
\end{frame}

% -----------------------------------------------------------------------------
% Slide 21
% -----------------------------------------------------------------------------
\begin{frame}[shrink=12]{Existential Example}
  ``Some robot is carrying a fragile package.''

  \[
    \exists x\exists y\;
    Robot(x) \land Package(y)
  \]
  \[
    \land\; Carries(x,y) \land Fragile(y).
  \]

  \begin{keyidea}
    The sentence claims that a suitable pair of objects exists.
  \end{keyidea}
\end{frame}

% -----------------------------------------------------------------------------
% Slide 22
% -----------------------------------------------------------------------------
\begin{frame}[shrink=10]{Quantifier Scope Changes Meaning}
  \begin{columns}[T,onlytextwidth]
    \begin{column}{0.47\textwidth}
      \[
        \forall x\exists y\; Loves(x,y)
      \]
      \begin{block}{Reading}
        Everyone loves someone.

        The ``someone'' may be different for each person.
      \end{block}
    \end{column}
    \begin{column}{0.47\textwidth}
      \[
        \exists y\forall x\; Loves(x,y)
      \]
      \begin{block}{Reading}
        There is one person whom everyone loves.

        The same $y$ works for all $x$.
      \end{block}
    \end{column}
  \end{columns}

  \begin{keyidea}
    Changing quantifier order can change the statement completely.
  \end{keyidea}
\end{frame}

% -----------------------------------------------------------------------------
% Slide 23
% -----------------------------------------------------------------------------
\begin{frame}[shrink=10]{A Classic Translation Trap}
  \begin{columns}[T,onlytextwidth]
    \begin{column}{0.54\textwidth}
      English often hides logical scope.

      \begin{block}{Sentence}
        Every student takes a course.
      \end{block}

      Does that mean each student may take a different course, or that there is one course every student takes?
    \end{column}
    \begin{column}{0.40\textwidth}
      \lectureimage[2.55in]{slide23.png}
    \end{column}
  \end{columns}
  \begin{keyidea}
    Formal logic is useful partly because it removes ambiguity that natural language tolerates.
  \end{keyidea}
\end{frame}

% -----------------------------------------------------------------------------
% Slide 24
% -----------------------------------------------------------------------------
\begin{frame}[shrink=12]{Translate: Every Student Takes at Least One Course}
  A faithful formalization is
  \[
    \forall x\; Student(x) \rightarrow
    \exists y\; (Course(y) \land Takes(x,y)).
  \]

  The existential quantifier sits \emph{inside} the universal.

  \begin{keyidea}
    Each student only needs some course; the course may be different for different students.
  \end{keyidea}
\end{frame}

% -----------------------------------------------------------------------------
% Slide 25
% -----------------------------------------------------------------------------
\begin{frame}[shrink=10]{Universal Statements Usually Use Implication}
  \begin{columns}[T,onlytextwidth]
    \begin{column}{0.47\textwidth}
      \[
        \forall x\; Student(x) \rightarrow HasID(x)
      \]
      \begin{block}{Why?}
        Objects that are not students are irrelevant to the rule.
      \end{block}
    \end{column}
    \begin{column}{0.47\textwidth}
      \[
        \forall x\; Student(x) \land HasID(x)
      \]
      \begin{block}{Problem}
        This says every object in the domain is both a student and has an ID.
      \end{block}
    \end{column}
  \end{columns}

  \begin{keyidea}
    For ``all $A$ are $B$,'' the common pattern is $\forall x\,(A(x)\rightarrow B(x))$.
  \end{keyidea}
\end{frame}

% -----------------------------------------------------------------------------
% Slide 26
% -----------------------------------------------------------------------------
\begin{frame}[shrink=10]{Existential Statements Usually Use Conjunction}
  \begin{columns}[T,onlytextwidth]
    \begin{column}{0.47\textwidth}
      \[
        \exists x\; Robot(x) \land Charging(x)
      \]
      \begin{block}{Reading}
        There exists an object that is both a robot and charging.
      \end{block}
    \end{column}
    \begin{column}{0.47\textwidth}
      \[
        \exists x\; Robot(x) \rightarrow Charging(x)
      \]
      \begin{block}{Problem}
        This can become true just by choosing an object that is not a robot.
      \end{block}
    \end{column}
  \end{columns}

  \begin{keyidea}
    For ``some $A$ is $B$,'' the common pattern is $\exists x\,(A(x)\land B(x))$.
  \end{keyidea}
\end{frame}

\section{Semantics and Knowledge Representation}

% -----------------------------------------------------------------------------
% Slide 27
% -----------------------------------------------------------------------------
\begin{frame}[shrink=10]{Semantics: FOL Still Needs a World}
  \begin{columns}[T,onlytextwidth]
    \begin{column}{0.54\textwidth}
      \begin{itemize}
        \item A \textbf{domain} specifies the objects we are talking about.
        \item An \textbf{interpretation} assigns meanings to constants, predicates, and functions.
        \item A sentence is true or false relative to that interpretation.
      \end{itemize}
    \end{column}
    \begin{column}{0.40\textwidth}
      \lectureimage[2.55in]{slide27.png}
    \end{column}
  \end{columns}
  \begin{keyidea}
    Syntax gives the form of a sentence; semantics connects the symbols to a possible world.
  \end{keyidea}
\end{frame}

% -----------------------------------------------------------------------------
% Slide 28
% -----------------------------------------------------------------------------
\begin{frame}[shrink=11]{One Symbolic KB, One Structured World}
  \begin{columns}[T,onlytextwidth]
    \begin{column}{0.31\textwidth}
      \begin{block}{Domain}
        \{Alice, Bob, Robot7, Room214, PackageA, $\ldots$\}
      \end{block}
    \end{column}
    \begin{column}{0.31\textwidth}
      \begin{block}{Interpretation}
        $Alice$ maps to a particular person.

        $In$ maps to a relation over ordered pairs.
      \end{block}
    \end{column}
    \begin{column}{0.31\textwidth}
      \begin{block}{Sentence}
        $In(Robot7,Room214)$ is true exactly when that pair belongs to the relation $In$.
      \end{block}
    \end{column}
  \end{columns}

  \lectureimage[1.20in]{slide28.png}

  \begin{keyidea}
    FOL separates the symbols we manipulate from the world those symbols are interpreted to describe.
  \end{keyidea}
\end{frame}

% -----------------------------------------------------------------------------
% Slide 29
% -----------------------------------------------------------------------------
\begin{frame}[shrink=10]{Knowledge Representation: Categories}
  \begin{columns}[T,onlytextwidth]
    \begin{column}{0.54\textwidth}
      Predicates naturally define categories such as
      \begin{itemize}
        \item $Student$;
        \item $Professor$;
        \item $Robot$;
        \item $Room$;
        \item $Package$.
      \end{itemize}

      Category relationships can then be expressed as general rules.
    \end{column}
    \begin{column}{0.40\textwidth}
      \lectureimage[2.60in]{slide29.png}
    \end{column}
  \end{columns}
  \begin{keyidea}
    Categories compress repeated facts by letting properties be inherited through rules.
  \end{keyidea}
\end{frame}

% -----------------------------------------------------------------------------
% Slide 30
% -----------------------------------------------------------------------------
\begin{frame}[shrink=10]{Taxonomies as Logical Rules}
  A simple hierarchy can be written as implications:
  \[
    \forall x\; GraduateStudent(x) \rightarrow Student(x)
  \]
  \[
    \forall x\; Student(x) \rightarrow Person(x)
  \]
  \[
    \forall x\; Person(x) \rightarrow Agent(x)
  \]

  \lectureimage[1.25in]{slide30.png}

  \begin{keyidea}
    A taxonomy organizes ``is-a'' relationships between categories.
  \end{keyidea}
\end{frame}

% -----------------------------------------------------------------------------
% Slide 31
% -----------------------------------------------------------------------------
\begin{frame}[shrink=12]{Inheritance Emerges from the Rules}
  Suppose the KB contains
  \[
    GraduateStudent(Maya)
  \]
  \[
    GraduateStudent(x) \rightarrow Student(x)
  \]
  \[
    Student(x) \rightarrow Person(x)
  \]
  \[
    Person(x) \rightarrow Agent(x).
  \]

  \begin{keyidea}
    The hierarchy makes higher-level properties available to lower-level categories without repeating them.
  \end{keyidea}
\end{frame}

% -----------------------------------------------------------------------------
% Slide 32
% -----------------------------------------------------------------------------
\begin{frame}[shrink=10]{Inheritance Is Powerful --- and Can Be Awkward}
  \begin{columns}[T,onlytextwidth]
    \begin{column}{0.54\textwidth}
      \begin{itemize}
        \item Simple logical inheritance is monotonic: adding facts does not retract earlier conclusions.
        \item Real categories often have exceptions.
        \item ``Birds fly'' is useful as a default, but penguins create an exception.
      \end{itemize}
    \end{column}
    \begin{column}{0.40\textwidth}
      \lectureimage[2.50in]{slide32.png}
    \end{column}
  \end{columns}
  \begin{keyidea}
    Classical FOL handles explicit exceptions, but default reasoning requires extra machinery.
  \end{keyidea}
\end{frame}

% -----------------------------------------------------------------------------
% Slide 33
% -----------------------------------------------------------------------------
\begin{frame}[shrink=10]{Relations Turn a Taxonomy into a Rich World Model}
  \begin{columns}[T,onlytextwidth]
    \begin{column}{0.38\textwidth}
      \[
        Student(Maya)
      \]
      \[
        EnrolledIn(Maya,AI)
      \]
      \[
        TaughtBy(AI,ProfLee)
      \]
      \[
        LocatedIn(AI,Building14)
      \]
    \end{column}
    \begin{column}{0.56\textwidth}
      \lectureimage[2.85in]{slide33.png}
    \end{column}
  \end{columns}
  \begin{keyidea}
    Unary predicates classify objects; multi-place predicates connect them into a network of knowledge.
  \end{keyidea}
\end{frame}

% -----------------------------------------------------------------------------
% Slide 34
% -----------------------------------------------------------------------------
\begin{frame}[shrink=10]{From Logical Relations to Knowledge Graphs}
  \begin{columns}[T,onlytextwidth]
    \begin{column}{0.54\textwidth}
      \begin{itemize}
        \item A graph view treats entities as nodes and relations as labeled edges.
        \item This is a convenient visualization of many FOL-style facts.
        \item But a graph alone does not automatically capture the full semantics of quantifiers and logical rules.
      \end{itemize}
    \end{column}
    \begin{column}{0.40\textwidth}
      \lectureimage[2.55in]{slide34.png}
    \end{column}
  \end{columns}
  \begin{keyidea}
    Knowledge graphs and logic are closely related representations, but they are not identical languages.
  \end{keyidea}
\end{frame}

% -----------------------------------------------------------------------------
% Slide 35
% -----------------------------------------------------------------------------
\begin{frame}[shrink=10]{Ontologies Add a Shared Vocabulary}
  \begin{columns}[T,onlytextwidth]
    \begin{column}{0.54\textwidth}
      An ontology specifies
      \begin{itemize}
        \item the kinds of entities that exist in a domain;
        \item important properties and relationships among those kinds;
        \item constraints or axioms that give the vocabulary precise meaning.
      \end{itemize}
    \end{column}
    \begin{column}{0.40\textwidth}
      \lectureimage[2.60in]{slide35.png}
    \end{column}
  \end{columns}
  \begin{keyidea}
    An ontology is a model of the domain's concepts, not merely a list of facts about individual objects.
  \end{keyidea}
\end{frame}

% -----------------------------------------------------------------------------
% Slide 36
% -----------------------------------------------------------------------------
\begin{frame}[shrink=10]{Campus Ontology Example}
  \begin{columns}[T,onlytextwidth]
    \begin{column}{0.30\textwidth}
      \begin{block}{Classes}
        Person

        Student

        Professor

        Course

        Room

        Building
      \end{block}
    \end{column}
    \begin{column}{0.34\textwidth}
      \begin{block}{Relations}
        $enrolledIn(Student,Course)$

        $teaches(Professor,Course)$

        $locatedIn(Course,Room)$

        $partOf(Room,Building)$
      \end{block}
    \end{column}
    \begin{column}{0.30\textwidth}
      \begin{block}{Rules / axioms}
        GraduateStudent $\subseteq$ Student

        Professor $\subseteq$ Person

        Every Course is located in some Room
      \end{block}
    \end{column}
  \end{columns}

  \begin{keyidea}
    An ontology gives many agents or systems a common conceptual schema for talking about the same domain.
  \end{keyidea}
\end{frame}

% -----------------------------------------------------------------------------
% Slide 37
% -----------------------------------------------------------------------------
\begin{frame}[shrink=10]{Taxonomy vs. Ontology}
  \begin{columns}[T,onlytextwidth]
    \begin{column}{0.47\textwidth}
      \begin{block}{Taxonomy}
        Primarily organizes categories into a hierarchy.

        GraduateStudent $\rightarrow$ Student $\rightarrow$ Person
      \end{block}
    \end{column}
    \begin{column}{0.47\textwidth}
      \begin{block}{Ontology}
        Includes categories plus richer relations, properties, constraints, and axioms about the domain.
      \end{block}
    \end{column}
  \end{columns}

  \begin{keyidea}
    A taxonomy answers ``what kind of thing is this?''; an ontology can also encode ``how can these kinds of things relate?''
  \end{keyidea}
\end{frame}

\section{Expressiveness and the Next Inference Problem}

% -----------------------------------------------------------------------------
% Slide 38
% -----------------------------------------------------------------------------
\begin{frame}[shrink=10]{Representation Choices Change the Inference Problem}
  \begin{columns}[T,onlytextwidth]
    \begin{column}{0.54\textwidth}
      \begin{itemize}
        \item A flat propositional KB may need thousands of repeated propositions.
        \item FOL can express the same regularities compactly with variables and quantifiers.
        \item That extra expressive power makes inference more complicated.
      \end{itemize}
    \end{column}
    \begin{column}{0.40\textwidth}
      \lectureimage[2.45in]{slide38.png}
    \end{column}
  \end{columns}
  \begin{keyidea}
    More expressive representations usually buy compactness at the cost of harder reasoning.
  \end{keyidea}
\end{frame}

% -----------------------------------------------------------------------------
% Slide 39
% -----------------------------------------------------------------------------
\begin{frame}[shrink=10]{Propositional Logic vs. First-Order Logic}
  \begin{center}
    \renewcommand{\arraystretch}{1.35}
    \begin{tabular}{l|p{0.31\textwidth}|p{0.37\textwidth}}
      & \textbf{Propositional logic} & \textbf{First-order logic} \\\hline
      Basic unit & Atomic proposition & Object + predicate / relation \\
      Variables & No object variables & Yes \\
      Quantifiers & No & $\forall$ and $\exists$ \\
      Compact general rules & Limited & Natural \\
      Inference & Simpler & More expressive, more complex \\
    \end{tabular}
  \end{center}

  \begin{keyidea}
    FOL is attractive because one symbolic pattern can cover arbitrarily many objects.
  \end{keyidea}
\end{frame}

% -----------------------------------------------------------------------------
% Slide 40
% -----------------------------------------------------------------------------
\begin{frame}[shrink=10]{What the Agent Can Now Represent}
  \begin{columns}[T,onlytextwidth]
    \begin{column}{0.54\textwidth}
      \begin{itemize}
        \item Specific facts: $In(Robot7,Room214)$.
        \item Categories: $Robot(Robot7)$.
        \item General rules: $\forall x\; Robot(x)\rightarrow Agent(x)$.
        \item Existence claims: $\exists x\; Charging(x)$.
        \item Relationships among arbitrary objects.
      \end{itemize}
    \end{column}
    \begin{column}{0.40\textwidth}
      \lectureimage[2.55in]{slide40.png}
    \end{column}
  \end{columns}
  \begin{keyidea}
    We now have a language expressive enough to describe structured domains rather than isolated Boolean facts.
  \end{keyidea}
\end{frame}

% -----------------------------------------------------------------------------
% Slide 41
% -----------------------------------------------------------------------------
\begin{frame}[shrink=10]{But How Does the Computer Use a Rule with Variables?}
  \begin{columns}[T,onlytextwidth]
    \begin{column}{0.54\textwidth}
      Suppose the KB contains
      \[
        \forall x\; Human(x)\rightarrow Mortal(x)
      \]
      and
      \[
        Human(Socrates).
      \]

      A human immediately sees that $x$ should correspond to $Socrates$.

      The inference engine needs an algorithm for making that match.
    \end{column}
    \begin{column}{0.40\textwidth}
      \lectureimage[2.65in]{slide41.png}
    \end{column}
  \end{columns}
  \begin{keyidea}
    The representation problem is solved; next we need machinery for matching variables to concrete terms.
  \end{keyidea}
\end{frame}

% -----------------------------------------------------------------------------
% Slide 42
% -----------------------------------------------------------------------------
\begin{frame}[shrink=12]{Next Time: Substitution and Unification}
  The key move will be finding a substitution such as
  \[
    \theta = \{x/Socrates\}.
  \]

  Applying it gives
  \[
    Human(x)\theta = Human(Socrates)
  \]
  and
  \[
    Mortal(x)\theta = Mortal(Socrates).
  \]

  \begin{keyidea}
    Unification turns general first-order rules into applicable instances.
  \end{keyidea}
\end{frame}

% -----------------------------------------------------------------------------
% Slide 43
% -----------------------------------------------------------------------------
\begin{frame}[shrink=10]{Takeaways}
  \begin{itemize}
    \item FOL represents objects, properties, relations, functions, and general rules.
    \item $\forall$ expresses universal claims; $\exists$ expresses existence claims; scope and order matter.
    \item Categories and taxonomies organize ``is-a'' structure and support inheritance.
    \item Ontologies add a richer shared model of domain concepts and relationships.
    \item Expressiveness reduces representational duplication but makes inference harder.
    \item Next: substitution, unification, and first-order inference.
  \end{itemize}

  \begin{keyidea}
    Knowledge representation chooses the language of thought; inference determines how an agent reasons with that language.
  \end{keyidea}
\end{frame}

\end{document}
